\subsection{Predicting Units with \unitytwo}\label{sec:offline:unity2}
Similar to \mfourt, the task of speech-to-speech translation in \mfourttwo is broken down into speech-to-text translation (\st) and then text-to-unit conversion (\tu).
While \unity relaxes the training difficulty of direct \sst, the \tu model often hallucinates or truncates the output.
This issue can be attributed to the length mismatch between the speech sequence in units and the text sequence in subwords, the former being on average 25 times longer.
To reduce errors in the unit generation, we propose a new direct two-pass {\sst} architecture, {\unitytwo}, which enhances the unit generation of {\unity} by a non-autoregressive (NAR) decoder that does not rely on any external aligner.

The overall architecture of \unitytwo is depicted in \Cref{fig:offline:unity2}.
\unitytwo replaces the second-pass autoregressive unit decoder in \unity with a NAR unit decoder.
We adopted the decoder architecture of FastSpeech2~\citep{ren2021fastspeech} and extended it to discrete unit generation.
The original FastSpeech2 decoder, designed for generating Mel-filterbank features in \tts, relies on a variance adaptor to add different variance information such as duration, pitch, and energy as conditional inputs before decoding.
Given that \unitytwo is designed to model discrete units, we only added duration information with a duration predictor; 
other information like pitch or prosody are not well-preserved by discrete units~\citep{polyak21_interspeech}.
The NAR unit decoder needs to expand the text input sequence to match the length of the unit output sequence, as such, a text-to-unit 
alignment is necessary.
Unlike \unity, \unitytwo predicts a duplicated (or non-reduced) unit sequence that preserves repetitive units.
Although the non-reduced unit sequence is longer, fast inference with NAR unit generation can compensate for the computational overhead.

\unitytwo starts with hierarchically upsampling of the \tu encoder output from subword-length, to character-length then to unit-length (\Cref{sec:offline:unity2:hie_up}).
The \textit{unit duration predictor}, key to the hierarchical upsampling, is supervised during training by a multilingual \textit{aligner} based on RAD-TTS~\citep{shih2021rad} (\Cref{sec:offline:unity2:unsup_aligner}).
To address the multimodality problem in the NAR generation of large-vocabulary non-reduced discrete units, we propose an efficient single-pass span-based glancing training (\Cref{sec:offline:glat}).




\subsubsection{Hierarchical subword-to-unit upsampling}\label{sec:offline:unity2:hie_up}
The \tu encoder in \unitytwo receives coarse subword-length representations from the \xt text decoder.\footnote{Subword vocabulary size of 256K in \mfourt}
As a first-pass decoder in \unitytwo, its features are not suitable for describing acoustic details necessary for subsequent unit prediction.
To leverage fine-grained textual information without hindering translation quality or the efficiency of the \xt decoder, 
we propose \textit{hierarchical subword-to-unit upsampling}, where we upsample the subword-length \tu encoder representations to character-length then to unit-length.

Specifically, a subword-to-character upsampler \texttt{Sub2Char} repeats each subword-length representation $\tuoutputstep$ according to the number of characters in the subword $\tgttextstep$ and adds character-level embeddings.
With $\tgttextchar$, the character-length sequence corresponding to $\tgttext$, 
we compute the character-length representations $\chartuoutput$ as follows:
\begin{align}\label{eq:char_align}
\duroutcharstep & = \subtwocharalignfunc(\tgttextstep, \tgttextchar), \\
j & = \sum_{l<i} \duroutcharstepl + m \ (m = 1 \cdots, \duroutcharstep), \\
\begin{split}
\chartuoutputstepj & = \texttt{Sub2Char}(\tuoutputstep, \tgttextcharstepj, j) \\
& = \tuoutputstep + \charemb(\tgttextcharstepj) + \frac{1}{\sqrt{d_{\text{model}}}} \cdot \charposemb(j),
\end{split}
\end{align}
where $\subtwocharalignfunc$ is a function returning the number of characters ($\duroutcharstep$) in the $i$-th subword,
$\charemb$ is a character embedding lookup table, 
$\charposemb$ is a character-level positional embedding layer,
and $d_{\text{model}}$ is the model dimension.

Then, a character-to-unit upsampler \texttt{Char2Unit} further upsamples $\chartuoutput$ to unit-length representations $\unittuoutput$ as:
\begin{align}
\duroutunitstep & = \chartwounitalignfunc(\tgttextcharstepj, \tgtunitnr, \hardalignment; \paradur), \label{eqn:offline:char2unit:dur} \\
k &= \sum_{l<j} \duroutunitstepl + m \ (m = 1, \cdots, \duroutunitstep), \\
\unittuoutputstepk & = \texttt{Char2Unit}(\chartuoutputstepj, k)  = \chartuoutputstepj + \alpha \cdot \unitposemb(k),
\end{align}
where $\chartwounitalignfunc$ is a duration predictor (parameterized by $\paradur$) that returns the number of duplicated units ($\duroutunitstep$) aligned with the $j$-th character,
$\hardalignment$ is a hard character-to-unit alignment matrix, 
$\alpha$ is a learnable scale parameter, 
and $\unitposemb$ is a unit-level positional embedding layer.
$\unittuoutput$ is used as input to the NAR unit decoder.
The duration predictor $\chartwounitalignfunc$ in \cref{eqn:offline:char2unit:dur} is trained to optimize $\lossdurpred(\paradur)$, a mean square error (MSE) loss taking the duration predicted by the aligner as training target in the logarithmic domain.

\subsubsection{Unsupervised multilingual character-to-unit alignment learning}\label{sec:offline:unity2:unsup_aligner}
For upsampling, the NAR unit decoder requires alignment ($\hardalignment$) between characters and units to train the unit duration predictor $\chartwounitalignfunc$.
The original FastSpeech2 used a forced alignment tool (\eg Montreal Forced Aligner~\citep{mcauliffe2017montreal}) to supervise the duration predictor.
For our massively multilingual efforts, forced aligners are unavailable for many low-resource languages.
To circumvent the need for external aligners, we propose an \textit{unsupervised multilingual character-to-unit aligner}.
We adapted the aligner architecture in RAD-TTS~\citep{shih2021rad} to our use case.
Namely, the \mfourttwo multilingual char-to-unit aligner is 
(1) modified to take discrete units and characters as inputs,
(2) trained in a multilingual fashion on \ntslangs languages, and
(3) trained with curriculum learning for the alignment prior.

For a character-length sequence $\tgttextchar$ and the associated unit sequence $\tgtunitnr$, let $s^\text{char}$ and $s^\text{unit}$ be the outputs of the aligner's two encoders (one for characters and one for units).
A soft alignment 
$\softalignment$ is calculated as follows:
\begin{align}
    \textrm{D}_{i,j} & = ||s^{\text{char}}_{i} - s^{\text{unit}}_{j}||_{2}, \\
    \softalignment_{i,j} & = \frac{e^{-\textrm{D}_{i,j}}}{\sum_{k}{e^{-D_{k,j}}}} + \textrm{P}_{\text{prior}}(i|j),
\end{align}
where $\textrm{P}_{\text{prior}}$ is the Beta-binomial alignment prior to encourage near-diagonal paths~\citep{shih2021rad}.
We disabled this alignment prior after 8k training steps to let the aligner learn a more accurate alignment later in the training.
To extract a hard alignment $\hardalignment$ from $\softalignment$, the monotonic alignment search (MAS) algorithm~\citep{kim2020glow} is applied. 

To optimize the aligner parameters $\paraaligner$, we maximized the log-likelihood of all possible monotonic alignment paths $\mathcal S(\tgttextchar)$, based on the forward algorithm.
The forward sum loss $\lossalignerfwd$ is formulated as:
\begin{align}\label{eq:fwd_sum_loss}
    \begin{split}
    \lossalignerfwd(\paraaligner) 
    & = - \log \textrm{P}\left(\mathcal S(\tgttextchar)|\tgtunitnr; \paraaligner\right), \\
    & = - \log \sum_{a \in \mathcal S(\tgttextchar)} \prod_{j=1}^{\unitlen} \textrm{P}(a_{j}|\tgtunitnr_{j}; \paraaligner),
    \end{split}
\end{align}
where the marginalization is efficiently implemented using a CTC loss.
To enforce that $\softalignment$ matches $\hardalignment$, a binarization loss $\lossalignerbin = \hardalignment \odot \log \softalignment$ with $\odot$ the Hadamard product.
This term is simply the KL divergence between the two alignments.
$\lossalignerbin$ is added after $K_{\text{bin}}$ training steps.

\subsubsection{Efficient span-based glancing training for NAR unit generation}\label{sec:offline:glat}
Non-autoregressive sequence generation suffers from the \textit{multimodality} problem.\footnote{Each token’s distribution depends only on the source sentence; this conditional independence assumption prevents a model from properly capturing the highly multimodal distribution of target translations~\citep{gu2017non}.}
Previous works have addressed this problem by using iterative decoding~\citep{lee-etal-2018-deterministic},
powerful generative models like normalizing flows~\citep{ma-etal-2019-flowseq}, 
or diffusion-based models~\citep{gong2022diffuseq,reid2022diffuser}.
In this work, we used a single-step NAR decoder to maintain inference efficiency.
Particularly, \unitytwo's NAR \tu decoder is a \textit{Glancing Transformer} (GLAT)~\citep{qian-etal-2021-glancing} that relaxes NAR token prediction by glancing at the ground-truth tokens.
When the naive GLAT based on random masking is used for unit prediction, the task becomes trivial since adjacent units are locally correlated.
To adapt GLAT to unit prediction, we propose an efficient \textit{span-based GLAT} that operates on the character length before glancing at the units.

Given a unit prediction accuracy $\alpha$, we sampled character positions to mask with a probability $1-\alpha$.
With $\mathcal I^{\text{char}}$ the set of the sampled positions to be masked in $\tgttextchar$,
we obtained the corresponding unit positions $\mathcal I^{\text{unit}}$ following the aligner's $\hardalignment$.
We then replaced the decoder input in the $\mathcal I^{\text{unit}}$ positions with ground-truth unit embeddings.
We demonstrate that the proposed span-based masking is more effective than random masking at the unit level.
Furthermore, we propose an efficient training based on a single forward pass where $\alpha$ is estimated from the previous $K_{\text{glat}}$ steps, instead of introducing a duplicate forward pass at each training step.

\subsubsection{Training \unitytwo's NAR \tu and aligner}
\label{sec:offline:unity2_training}

The second-pass NAR unit decoder and aligner are jointly trained with the following objective:
\begin{align}\label{eq:tu_total_loss}
\begin{split}
\lossnartutotal(\paratu, \paradur, \paraaligner)
&= 
\lossnartuce(\paratu) + 
\lossdurpred(\paradur) +
\losscharinterctc(\paratu) 
\\ &\phantom{=}+
\lossalignerfwd(\paraaligner) + 
\lossalignerbin(\paraaligner), 
\end{split}
\end{align}
where $\losscharinterctc$ is a character-level CTC loss at an intermediate unit decoder layer, added to accelerate training convergence~\citep{lee2021intermediate}.

